\chapter{Wprowadzenie i analiza problemu}
\label{ch:wprowadzenie}

\section{Cel i zakres projektu}

Celem projektu by\l{}o stworzenie aplikacji desktopowej wspomagaj\k{a}cej
prac\k{e} o\'{s}rodka szkolenia kierowc\'{o}w. Tradycyjne OSK prowadz\k{a}
dokumentacj\k{e} w~formie papierowej lub w~arkuszach kalkulacyjnych, co jest
czasoch\l{}onne i~podatne na b\l{}\k{e}dy. Proponowana aplikacja automatyzuje
najwa\.{z}niejsze czynno\'{s}ci -- zarz\k{a}dzanie kursantami, instruktorami,
pojazdami i~rezerwacjami jazd.

Zakres projektu obejmuje:
\begin{itemize}
  \item aplikacj\k{e} okienkow\k{a} dla systemu Windows,
  \item relacyjn\k{a} baz\k{e} danych obs\l{}ugiwan\k{a} przez kod C\#,
  \item system logowania z~trzema rolami u\.{z}ytkownik\'{o}w,
  \item kalendarz rezerwacji termin\'{o}w jazd,
  \item walidacj\k{e} wprowadzanych danych (m.in. numer PESEL),
  \item bezpieczne przechowywanie hase\l{} i~ochron\k{e} danych.
\end{itemize}

\section{Przegl\k{a}d istniej\k{a}cych rozwi\k{a}za\'{n}}
\label{sec:przeglad}

Na rynku dost\k{e}pnych jest kilka aplikacji przeznaczonych dla o\'{s}rodk\'{o}w
szkolenia kierowc\'{o}w. Poni\.{z}ej om\'{o}wiono trzy przyk\l{}ady.

\textbf{iMSZ (Internetowy Modu\l{} Szko\l{}y Jazdy)} to popularne w~Polsce
rozwi\k{a}zanie webowe oferuj\k{a}ce zarz\k{a}dzanie kursantami i~harmonogramem
jazd. Dzia\l{}a w~przegl\k{a}darce, lecz wymaga sta\l{}ego dost\k{e}pu do
internetu i~op\l{}aty miesi\k{e}cznej.

\textbf{OSK Zintegrowany} to aplikacja desktopowa dla Windows z~modu\l{}em
kalendarza. Jest to oprogramowanie p\l{}atne i~nie wspiera nowoczesnych
standard\'{o}w bezpiecze\'{n}stwa.

\textbf{Proton OSK} to r\'{o}wnie\.{z} aplikacja okienkowa z~modu\l{}em
rezerwacji i~raportami. Wymaga oddzielnej licencji i~nie oferuje
bezp\l{}atnego wdro\.{z}enia.

Projekt realizowany w~ramach niniejszej pracy stanowi uproszczon\k{a}, ale
w~pe\l{}ni dzia\l{}aj\k{a}c\k{a} alternatyw\k{e} bez zewn\k{e}trznych koszt\'{o}w.
